\chapter{Fondamenti di Algebra Lineare Numerica e Decomposizione SVD}
\label{chap:fondamenti-algebra-lineare-svd}

\FloatBarrier
\section{Il Ruolo dell'Algebra Lineare nel Machine Learning e nella Data Analysis}

Il moderno apprendimento automatico e l'analisi su larga scala dei dati poggiano in larga misura su concetti e strumenti derivati dall'algebra lineare numerica e dall'ottimizzazione continua~\cite{strang2019linear,golub2013matrix}. I dati del mondo reale (immagini, testi, serie temporali, grafi o segnali sensoriali) vengono sistematicamente codificati sotto forma di vettori e matrici numeriche:

\begin{itemize}
    \item \textbf{Matrice dei dati (Design Matrix):} Sia $\mathcal{D} = \{\mathbf{x}_1, \dots, \mathbf{x}_n\}$ un insieme di $n$ campioni statistici, ciascuno caratterizzato da $d$ feature numeriche. L'intera collezione si rappresenta come matrice $X \in \R^{n \times d}$:
    \begin{equation}
        X = \begin{bmatrix}
            \mathbf{x}_1^T \\
            \mathbf{x}_2^T \\
            \vdots \\
            \mathbf{x}_n^T
        \end{bmatrix} = \begin{bmatrix}
            x_{11} & x_{12} & \cdots & x_{1d} \\
            x_{21} & x_{22} & \cdots & x_{2d} \\
            \vdots & \vdots & \ddots & \vdots \\
            x_{n1} & x_{n2} & \cdots & x_{nd}
        \end{bmatrix},
    \end{equation}
    dove le righe corrispondono alle osservazioni e le colonne alle variabili misurate.
    \item \textbf{Matrice di Gram e Matrice di Covarianza:} Le relazioni lineari tra coppie di osservazioni o coppie di variabili si quantificano mediante prodotti matriciali: la matrice di Gram $G = X X^T \in \R^{n \times n}$ quantifica le similarità tra campioni, mentre la matrice di covarianza campionaria non centrata $\Sigma_X = \frac{1}{n} X^T X \in \R^{d \times d}$ misura la varianza e covarianza congiunta tra le feature.
\end{itemize}

A differenza dell'algebra lineare simbolica o teorica, l'algebra lineare \emph{computazionale} o \emph{numerica} deve affrontare due vincoli ineludibili: il costo computazionale asintotico delle operazioni su larga scala e la propagazione degli errori dovuti all'aritmetica di macchina a precisione finita (standard IEEE 754)~\cite{trefethen1997numerical}.

\FloatBarrier
\section{Norme Vettoriali e Matriciali}

La quantificazione della magnitudine di vettori e matrici, nonché la misura di errore e distanza tra predizioni ed etichette reali, richiede la definizione assiomatica di \textbf{norma}.

\begin{definition}[Norma Vettoriale]
Una funzione $\|\cdot\|: \R^d \to \R$ si definisce \textbf{norma} su $\R^d$ se soddisfa le seguenti tre proprietà per ogni $\mathbf{x}, \mathbf{y} \in \R^d$ e $\alpha \in \R$:
\begin{enumerate}
    \item \textbf{Non-negatività e positività:} $\|\mathbf{x}\| \ge 0$, con $\|\mathbf{x}\| = 0 \iff \mathbf{x} = \mathbf{0}$;
    \item \textbf{Omogeneità assoluta:} $\|\alpha \mathbf{x}\| = |\alpha| \, \|\mathbf{x}\|$;
    \item \textbf{Disuguaglianza triangolare:} $\|\mathbf{x} + \mathbf{y}\| \le \|\mathbf{x}\| + \|\mathbf{y}\|$.
\end{enumerate}
\end{definition}

\subsection{Famiglia delle \texorpdfstring{$\ell_p$}{lp}-Norme}

Le norme vettoriali più utilizzate nell'analisi dei dati appartengono alla famiglia delle norme $\ell_p$ (con $p \ge 1$):
\begin{equation}
    \|\mathbf{x}\|_p = \left( \sum_{i=1}^d |x_i|^p \right)^{1/p}.
\end{equation}

I tre casi fondamentali sono:
\begin{enumerate}
    \item \textbf{Norma di Manhattan ($\ell_1$):} $\|\mathbf{x}\|_1 = \sum_{i=1}^d |x_i|$. Fondamentale per indurre sparsità nei modelli lineari (regolarizzazione Lasso).
    \item \textbf{Norma Euclidea ($\ell_2$):} $\|\mathbf{x}\|_2 = \sqrt{\sum_{i=1}^d x_i^2} = \sqrt{\mathbf{x}^T \mathbf{x}}$. Rappresenta la distanza geometrica ordinaria ed è invariante per rotazioni ortogonali.
    \item \textbf{Norma di Chebyshev ($\ell_\infty$):} $\|\mathbf{x}\|_\infty = \max_{1 \le i \le d} |x_i|$. Misura lo scostamento massimo su singola coordinata.
\end{enumerate}

\begin{property}[Equivalenza Topologica delle Norme]
In uno spazio vettoriale a dimensione finita $\R^d$, tutte le norme sono topologicamente equivalenti: per ogni coppia di norme $\|\cdot\|_a$ e $\|\cdot\|_b$, esistono costanti positive $c_1, c_2 > 0$ tali che:
\begin{equation}
    c_1 \|\mathbf{x}\|_b \le \|\mathbf{x}\|_a \le c_2 \|\mathbf{x}\|_b \qquad \forall \mathbf{x} \in \R^d.
\end{equation}
\end{property}

\subsection{Norme Matriciali e Norme Subordinate}

Per quantificare la dimensione di una trasformazione lineare $A \in \R^{m \times n}$ si utilizzano norme matriciali. Si distinguono due grandi categorie:

\paragraph{Norma di Frobenius.}
Estende direttamente la norma Euclidea trattando la matrice come un lungo vettore di dimensione $m \cdot n$:
\begin{equation}
    \|A\|_F = \sqrt{\sum_{i=1}^m \sum_{j=1}^n a_{ij}^2} = \sqrt{\trace(A^T A)}.
\end{equation}

\paragraph{Norme Indotte (o Subordinate).}
Misurano la massima amplificazione che la matrice $A$ può imprimere a un vettore non nullo rispetto a una norma vettoriale data:
\begin{equation}
    \|A\|_p = \sup_{\mathbf{x} \ne \mathbf{0}} \frac{\|A \mathbf{x}\|_p}{\|\mathbf{x}\|_p} = \max_{\|\mathbf{x}\|_p = 1} \|A \mathbf{x}\|_p.
\end{equation}

Le espressioni computazionali notevoli sono:
\begin{itemize}
    \item $\|A\|_1 = \max_{1 \le j \le n} \sum_{i=1}^m |a_{ij}|$ (massima somma assoluta sulle colonne);
    \item $\|A\|_\infty = \max_{1 \le i \le m} \sum_{j=1}^n |a_{ij}|$ (massima somma assoluta sulle righe);
    \item $\|A\|_2 = \sigma_{\max}(A) = \sqrt{\lambda_{\max}(A^T A)}$ (\textbf{norma spettrale}, pari al massimo valore singolare di $A$).
\end{itemize}

Ogni norma indotta è \textbf{sub-moltiplicativa}: $\|A B\| \le \|A\| \, \|B\|$.

\FloatBarrier
\section{Condizionamento e Stabilità Numerica}

Un problema computazionale $y = f(x)$ si definisce \textbf{ben condizionato} se a piccole perturbazioni sui dati di ingresso $x$ corrispondono piccole variazioni sul risultato $y$. Al contrario, un problema si definisce \textbf{mal condizionato} (\textit{ill-conditioned}) se minime perturbazioni nei dati producono oscillazioni macroscopicamente ampie nella soluzione~\cite{golub2013matrix}.

\subsection{Numero di Condizionamento di un Sistema Lineare}

Si consideri il sistema lineare quadrato $A \mathbf{x} = \mathbf{b}$, con $A \in \R^{n \times n}$ non singolare e $\mathbf{b} \ne \mathbf{0}$. Si supponga che il termine noto sia affetto da un rumore di misura $\delta \mathbf{b}$, generando una variazione $\delta \mathbf{x}$ nella soluzione:
\begin{equation}
    A(\mathbf{x} + \delta \mathbf{x}) = \mathbf{b} + \delta \mathbf{b} \implies A \delta \mathbf{x} = \delta \mathbf{b} \implies \delta \mathbf{x} = A^{-1} \delta \mathbf{b}.
\end{equation}

Dalla sub-moltiplicatività della norma segue:
\begin{equation}
    \|\delta \mathbf{x}\| \le \|A^{-1}\| \, \|\delta \mathbf{b}\| \qquad \text{e} \qquad \|\mathbf{b}\| \le \|A\| \, \|\mathbf{x}\| \implies \frac{1}{\|\mathbf{x}\|} \le \frac{\|A\|}{\|\mathbf{b}\|}.
\end{equation}

Moltiplicando le due disuguaglianze membro a membro si ottiene la relazione fondamentale:
\begin{equation}
    \frac{\|\delta \mathbf{x}\|}{\|\mathbf{x}\|} \le \left(\|A\| \, \|A^{-1}\|\right) \frac{\|\delta \mathbf{b}\|}{\|\mathbf{b}\|}.
\end{equation}

\begin{definition}[Numero di Condizionamento]
Il \textbf{numero di condizionamento} di una matrice quadrata invertibile $A \in \R^{n \times n}$ associato alla norma $\|\cdot\|$ è il valore:
\begin{equation}
    \kappa(A) = \|A\| \, \|A^{-1}\|.
\end{equation}
Nel caso della norma spettrale $\|\cdot\|_2$, indicando con $\sigma_1 \ge \sigma_2 \ge \dots \ge \sigma_n > 0$ i valori singolari di $A$:
\begin{equation}
    \kappa_2(A) = \|A\|_2 \, \|A^{-1}\|_2 = \frac{\sigma_{\max}(A)}{\sigma_{\min}(A)} = \frac{\sigma_1}{\sigma_n} \ge 1.
\end{equation}
\end{definition}

\begin{alertblock}{Significato Computazionale di $\kappa(A)$}
In aritmetica floating-point a doppia precisione (con precisione di macchina $\epsilon_{\text{mach}} \approx 10^{-16}$), la risoluzione di un sistema con $\kappa(A) \approx 10^k$ comporta tipicamente la perdita di circa $k$ cifre decimali significative. Se $\kappa(A) \ge 10^{16}$, la matrice è considerata numericamente singolare.
\end{alertblock}

\FloatBarrier
\section{Decomposizione Spettrale (Eigendecomposition)}

Sia $A \in \R^{n \times n}$ una matrice quadrata. Uno scalare $\lambda \in \C$ è un \textbf{autovalore} di $A$ con corrispondente \textbf{autovettore} non nullo $\mathbf{v} \ne \mathbf{0}$ se soddisfa l'equazione agli autovalori:
\begin{equation}
    A \mathbf{v} = \lambda \mathbf{v}.
\end{equation}

\begin{theorem}[Teorema Spettrale per Matrici Simmetriche Reali]
Sia $A \in \R^{n \times n}$ una matrice simmetrica reale ($A = A^T$). Allora:
\begin{enumerate}
    \item Tutti gli autovalori di $A$ sono reali: $\lambda_1, \dots, \lambda_n \in \R$.
    \item Gli autovettori associati ad autovalori distinti sono mutuamente ortogonali.
    \item Esiste una matrice ortogonale $Q \in \R^{n \times n}$ ($Q^T Q = Q Q^T = I_n$) le cui colonne sono autovettori normalizzati di $A$, tale che:
    \begin{equation}
        A = Q \Lambda Q^T = \sum_{i=1}^n \lambda_i \mathbf{q}_i \mathbf{q}_i^T,
    \end{equation}
    dove $\Lambda = \diag(\lambda_1, \dots, \lambda_n)$.
\end{enumerate}
\end{theorem}

\subsection{Matrici Semidefinite Positive (SPSD) e Forme Quadratiche}

Una matrice simmetrica $A \in \R^{n \times n}$ si definisce:
\begin{itemize}
    \item \textbf{Semidefinita positiva} ($A \succeq 0$) se $\mathbf{x}^T A \mathbf{x} \ge 0$ per ogni $\mathbf{x} \in \R^n \iff \lambda_i(A) \ge 0$ per tutti gli $i=1,\dots,n$.
    \item \textbf{Definita positiva} ($A \succ 0$) se $\mathbf{x}^T A \mathbf{x} > 0$ per ogni $\mathbf{x} \ne \mathbf{0} \iff \lambda_i(A) > 0$ per tutti gli $i=1,\dots,n$.
\end{itemize}

\begin{example}[La Matrice di Covarianza è sempre Semidefinita Positiva]
Sia $X \in \R^{n \times d}$ una matrice centrata con media campionaria nulla. La matrice di covarianza $C = \frac{1}{n} X^T X$ è simmetrica e per ogni vettore $\mathbf{w} \in \R^d$ vale:
\begin{equation}
    \mathbf{w}^T C \mathbf{w} = \mathbf{w}^T \left(\frac{1}{n} X^T X\right) \mathbf{w} = \frac{1}{n} (X \mathbf{w})^T (X \mathbf{w}) = \frac{1}{n} \|X \mathbf{w}\|_2^2 \ge 0.
\end{equation}
Dunque tutti gli autovalori di $C$ sono reali e non negativi, confermando che $C \succeq 0$.
\end{example}

\FloatBarrier
\section{Decomposizione ai Valori Singolari (SVD)}
\label{sec:svd}

La decomposizione spettrale è applicabile esclusivamente a matrici quadrate simmetriche. La \textbf{Singular Value Decomposition (SVD)} costituisce la generalizzazione fondamentale dell'algebra lineare a matrici rettangolari arbitrarie $A \in \R^{m \times n}$~\cite{golub2013matrix,strang2019linear}.

\begin{theorem}[Esistenza della SVD]
Data una qualsiasi matrice reale $A \in \R^{m \times n}$ di rango $r \le \min(m, n)$, esistono due matrici ortogonali $U \in \R^{m \times m}$ e $V \in \R^{n \times n}$ e una matrice pseudo-diagonale $\Sigma \in \R^{m \times n}$ tali che:
\begin{equation}
    A = U \Sigma V^T,
\end{equation}
dove $\Sigma$ presenta sulla diagonale principale i valori singolari disposti in ordine decrescente:
\begin{equation}
    \sigma_1 \ge \sigma_2 \ge \dots \ge \sigma_r > \sigma_{r+1} = \dots = \sigma_{\min(m, n)} = 0.
\end{equation}
\end{theorem}

\begin{figure}[H]
\centering
\begin{tikzpicture}[
    boxU/.style={rectangle, draw=blue!80!black, fill=blue!10, thick, minimum width=1.8cm, minimum height=2.8cm, align=center, font=\small\bfseries},
    boxS/.style={rectangle, draw=green!70!black, fill=green!10, thick, minimum width=2.4cm, minimum height=2.8cm, align=center, font=\small\bfseries},
    boxV/.style={rectangle, draw=orange!80!black, fill=orange!10, thick, minimum width=2.4cm, minimum height=1.6cm, align=center, font=\small\bfseries},
    boxA/.style={rectangle, draw=purple!80!black, fill=purple!10, thick, minimum width=2.4cm, minimum height=2.8cm, align=center, font=\small\bfseries},
    lbl/.style={font=\small}
]
    \node[boxA] (A) at (0, 0) {$A$\\$(m \times n)$};
    \node[lbl] at (1.8, 0) {$=$};
    \node[boxU] (U) at (3.2, 0) {$U$\\$(m \times m)$};
    \node[boxS] (S) at (5.7, 0) {$\Sigma$\\$(m \times n)$};
    \node[boxV] (V) at (8.5, 0.6) {$V^T$\\$(n \times n)$};

    \node[lbl, below=0.2cm of A, text=purple!80!black] {Dati originali};
    \node[lbl, below=0.2cm of U, text=blue!80!black] {Vettori singolari sx};
    \node[lbl, below=0.2cm of S, text=green!70!black] {Valori singolari};
    \node[lbl, below=1.4cm of V, text=orange!80!black] {Vettori singolari dx};
\end{tikzpicture}
\caption{Rappresentazione a blocchi dimensionali della Decomposizione ai Valori Singolari (SVD completa).}
\label{fig:svd-blocks}
\end{figure}

\subsection{Formulazione Esterna e SVD Compatta}

La matrice $A$ può essere espressa come somma di $r$ matrici di rango $1$:
\begin{equation}
    A = \sum_{i=1}^r \sigma_i \mathbf{u}_i \mathbf{v}_i^T,
\end{equation}
dove:
\begin{itemize}
    \item $\mathbf{u}_i \in \R^m$ è la $i$-esima colonna di $U$ (\textbf{vettore singolare sinistro}), autovettore di $A A^T$ con autovalore $\lambda_i = \sigma_i^2$;
    \item $\mathbf{v}_i \in \R^n$ è la $i$-esima colonna di $V$ (\textbf{vettore singolare destro}), autovettore di $A^T A$ con autovalore $\lambda_i = \sigma_i^2$;
    \item $\sigma_i = \sqrt{\lambda_i(A^T A)} = \sqrt{\lambda_i(A A^T)}$ è l'$i$-esimo valore singolare.
\end{itemize}

Nella pratica computazionale, se $m \gg n$, si adotta la \textbf{thin SVD} (o SVD economica), trattenendo unicamente le prime $n$ colonne di $U$ e la sottomatrice quadrata $\Sigma_n \in \R^{n \times n}$, con risparmio lineare di memoria $\BigO{m n}$ anziché $\BigO{m^2}$.

\FloatBarrier
\section{Proprietà ed Applicazioni Fondamentali di SVD}

\subsection{Miglior Approssimazione di Basso Rango: Teorema di Eckart-Young-Mirsky}

Nel data mining, nella compressione di segnali e nella riduzione della dimensionalità, si ricerca frequentemente una matrice $\hat{A}_k$ di rango ristretto $k < r$ che approssimi al meglio la matrice originaria $A$.

\begin{theorem}[Teorema di Eckart-Young-Mirsky]
Sia $A = \sum_{i=1}^r \sigma_i \mathbf{u}_i \mathbf{v}_i^T$ la SVD di $A \in \R^{m \times n}$. Per qualsiasi intero $1 \le k < r$, la matrice di rango al più $k$ definita dalla SVD troncata:
\begin{equation}
    A_k = \sum_{i=1}^k \sigma_i \mathbf{u}_i \mathbf{v}_i^T = U_k \Sigma_k V_k^T
\end{equation}
è la soluzione ottima del problema di miglior approssimazione sia in norma spettrale sia in norma di Frobenius:
\begin{equation}
    A_k = \argmin_{B \in \R^{m \times n},\, \rank(B) \le k} \|A - B\|_2 = \argmin_{B \in \R^{m \times n},\, \rank(B) \le k} \|A - B\|_F.
\end{equation}
Inoltre, gli errori minimi ottimali valgono esattamente:
\begin{align}
    \|A - A_k\|_2 &= \sigma_{k+1}, \\
    \|A - A_k\|_F &= \sqrt{\sum_{j=k+1}^r \sigma_j^2}.
\end{align}
\end{theorem}

\subsection{Pseudoinversa di Moore-Penrose e Minimi Quadrati}

Quando una matrice rettangolare $A \in \R^{m \times n}$ non è invertibile ($m \ne n$ o rango carente), la SVD fornisce lo strumento naturale per definire la \textbf{pseudoinversa di Moore-Penrose} $A^\dagger \in \R^{n \times m}$:
\begin{equation}
    A^\dagger = V \Sigma^\dagger U^T = \sum_{i=1}^r \frac{1}{\sigma_i} \mathbf{v}_i \mathbf{u}_i^T,
\end{equation}
dove $\Sigma^\dagger$ si ottiene trasponendo $\Sigma$ e invertendo i soli elementi diagonali non nulli $\sigma_i > 0$.

\begin{property}[Proprietà di Minimi Quadrati]
Dato il problema dei minimi quadrati lineari $\min_{\mathbf{x} \in \R^n} \|A \mathbf{x} - \mathbf{b}\|_2$, il vettore
\begin{equation}
    \mathbf{x}^* = A^\dagger \mathbf{b}
\end{equation}
è l'unica soluzione a norma minima $\|\mathbf{x}^*\|_2$ tra tutte le soluzioni che minimizzano il residuo quadratico.
\end{property}

\subsection{Connessione tra SVD e Principal Component Analysis (PCA)}

Sia $X \in \R^{n \times d}$ una matrice di $n$ dati centrati per colonna (media campionaria $\boldsymbol{\mu} = \mathbf{0}$). La matrice di covarianza campionaria vale:
\begin{equation}
    C = \frac{1}{n-1} X^T X.
\end{equation}
Calcolando la SVD di $X$: $X = U \Sigma V^T$, si ricava:
\begin{equation}
    C = \frac{1}{n-1} (V \Sigma U^T) (U \Sigma V^T) = \frac{1}{n-1} V \Sigma^2 V^T = V \left( \frac{\Sigma^2}{n-1} \right) V^T.
\end{equation}

Da questa uguaglianza discendono immediatamente tre conclusioni cruciali:
\begin{enumerate}
    \item I vettori singolari destri $\mathbf{v}_i$ di $X$ (colonne di $V$) sono esattamente gli \textbf{autovettori della matrice di covarianza} (direzioni principali della PCA).
    \item La varianza campionaria spiegata dalla $i$-esima componente principale è legata al valore singolare dalla relazione $\lambda_i = \frac{\sigma_i^2}{n-1}$.
    \item Le proiezioni dei dati (coordinate dei punteggi / \textit{scores}) nello spazio ridotto a $k$ dimensioni si calcolano senza mai costruire esplicitamente la matrice $C \in \R^{d \times d}$:
    \begin{equation}
        Z_k = X V_k = (U_k \Sigma_k V_k^T) V_k = U_k \Sigma_k.
    \end{equation}
\end{enumerate}

\begin{lstlisting}[style=pythoncode,caption={Calcolo compatto di Truncated SVD e PCA con NumPy.}]
import numpy as np

def compute_svd_pca(X: np.ndarray, k: int):
    """
    Calcola la proiezione PCA a k dimensioni sfruttando la SVD.
    X: matrice dei dati (n_samples, n_features)
    k: numero di componenti principali da trattenere
    """
    # 1. Centratura dei dati (media zero su ogni colonna)
    X_centered = X - np.mean(X, axis=0)
    n = X.shape[0]

    # 2. Decomposizione SVD economica
    U, s, Vt = np.linalg.svd(X_centered, full_matrices=False)

    # 3. Troncamento ai primi k fattori
    Uk = U[:, :k]
    sk = s[:k]
    Vk = Vt[:k, :].T

    # 4. Calcolo delle proiezioni (scores) e varianza spiegata
    scores = Uk * sk  # equivalente a X_centered @ Vk
    explained_variance = (sk ** 2) / (n - 1)
    explained_variance_ratio = explained_variance / np.sum((s ** 2) / (n - 1))

    return scores, Vk, explained_variance_ratio
\end{lstlisting}

\FloatBarrier
\section{Sintesi dei Concetti Chiave}

\begin{table}[H]
  \centering
  \small
  \setlength{\tabcolsep}{6pt}
  \begin{tabularx}{\textwidth}{l Y Y}
    \toprule
    \textbf{Strumento Matematico} & \textbf{Definizione Operativa} & \textbf{Impatto nella Data Analysis / ML} \\
    \midrule
    Norma $\ell_2$ e $\|A\|_F$ & Radice della somma dei quadrati & Ottimizzazione least squares, misura di errore MSE. \\
    Norma $\ell_1$ & Somma dei valori assoluti & Sparsità, feature selection robusta (Lasso). \\
    Numero di condizionamento $\kappa$ & $\sigma_{\max} / \sigma_{\min}$ & Stabilità numerica, amplificazione del rumore nei dati. \\
    Eigendecomposition & $A = Q \Lambda Q^T$ (solo per $A = A^T$) & Analisi spettrale, forme quadratiche, kernel positivi. \\
    SVD completa & $A = U \Sigma V^T$ (qualsiasi rettangolare) & Decomposizione geometrica universale dei dati. \\
    Eckart-Young-Mirsky & $A_k = \sum_{i=1}^k \sigma_i \mathbf{u}_i \mathbf{v}_i^T$ & Compressione ottimale a basso rango, filtro di rumore. \\
    SVD per PCA & $Z_k = U_k \Sigma_k$ & Riduzione dimensionale senza formare la matrice $X^T X$. \\
    \bottomrule
  \end{tabularx}
  \caption{Quadro riassuntivo dei concetti di algebra lineare numerica e SVD introdotti nel capitolo.}
  \label{tab:sintesi-algebra-lineare-svd}
\end{table}
